\chapter*{Заключение}
\addcontentsline{toc}{chapter}{Заключение}
\label{ch:conclusion}

В ходе выполнения лабораторной работы~\cfgLabNumber{} был создан макет
страницы интернет-магазина домашних растений в графическом редакторе
\textbf{Figma} в трёх адаптивных версиях: десктоп (\(1440\,\)px),
планшет (\(768\,\)px) и мобильный телефон (\(360\,\)px).

Освоенные инструменты и приёмы:
\begin{itemize}
  \item создание и настройка фреймов (\texttt{Frame}) как контейнеров
        макета;
  \item работа с текстовыми слоями: подбор шрифта, кегля, начертания,
        выравнивания;
  \item построение геометрических объектов и работа с цветом
        (\texttt{Fill}, \texttt{Stroke}, \texttt{Corner radius},
        тени);
  \item импорт фотографий через плагин \texttt{Unsplash};
  \item группировка элементов карточки товара во фрейм для последующего
        тиражирования (\texttt{Ctrl+Alt+G});
  \item настройка \texttt{Constraints} (привязок) для адаптивного
        поведения элементов;
  \item построение колоночной \texttt{Layout grid} как направляющей
        сетки;
  \item подготовка нескольких адаптивных версий копированием фрейма и
        изменением его ширины.
\end{itemize}

Полученные навыки --- основа повседневной работы веб-дизайнера. Концепции
фрейма, привязок и сетки имеют прямые аналоги в HTML/CSS
(\texttt{<div>}, Flexbox/Grid, медиа-запросы), что облегчает
последующий переход от макета к реальной вёрстке страницы.

\endinput
